\begin{abstract}
Large Language Models (LLMs) have demonstrated remarkable capabilities on
text-to-SQL benchmarks such as Spider and BIRD. However, these benchmarks
evaluate \emph{query generation}, not \emph{structural understanding} of
database schemas. We introduce \textbf{DW-Bench}, the first benchmark
specifically designed to evaluate LLMs on \emph{graph topology reasoning}
over data warehouse schemas. DW-Bench contains \textbf{674 questions} across
\textbf{4 real-world datasets} (AdventureWorks, TPC-DS, TPC-DI, OMOP~CDM)
spanning \textbf{198 tables}, \textbf{13 question subtypes}, and \textbf{3
difficulty levels}. Questions require reasoning about foreign key paths,
data lineage chains, transitive impact analysis, and connected component
detection---tasks that demand graph traversal rather than SQL generation.

We evaluate three baselines---\emph{Flat Text}, \emph{Vector-RAG}, and
\emph{Graph-Augmented}---using two frontier models: Gemini~2.5 Flash
(closed-weight) and DeepSeek-V3 (open-weight, 671B MoE). For Gemini,
micro-averaging creates a \textbf{Triviality Illusion}: Flat Text appears
strongest (73.9\% EM) by inflating scores through 316 easy routing
questions, but under macro-averaging Graph-Augmented leads
(\textbf{68.0\%} vs.\ 61.1\%). For DeepSeek-V3, Graph-Augmented
\emph{strictly dominates} across both micro (\textbf{82.3\%} vs.\
68.4\%) and macro (\textbf{73.9\%} vs.\ 56.1\%) metrics---confirming
that topological context injection is universally beneficial. An Oracle
baseline injecting gold graph evidence achieves \textbf{97--99\% EM},
proving these questions are solvable and isolating the failure as a
retrieval gap, not question ambiguity.

DW-Bench includes an obfuscation protocol replacing table names with
random identifiers. DeepSeek suffers up to \textbf{30\%} penalty,
revealing semantic reliance, while Gemini's hard task performance
paradoxically \emph{improves} (+24\%), suggesting that shorter random
identifiers reduce multi-hop tracking complexity.
\end{abstract}
